\section{总结与延伸}

\subsection{核心结论}

讲者经过整个讲座的论证，得出了一个清晰的结论：

\begin{importantbox}{{Your AI, Your Responsibility}}
\textbf{你的 AI 不能替你承担责任}。不要指望 AI 系统自己``做正确的事''——它没有道德判断力，没有因果推理能力，也没有真正的``理解''。如果你构建并发布了一个 AI 系统，那么这个系统的所有行为——无论是正确的还是有害的——都是\textbf{你、你的团队、你的公司}的责任。希波克拉底誓言应该由 AI 的\textbf{创造者}来宣誓，而不是 AI 本身。
\end{importantbox}

\subsection{从讲座中提炼的实践原则}

综合整个讲座的内容，可以提炼出以下实践原则：

\begin{enumerate}
    \item \textbf{评估风险-收益}：在启动 AI 项目前，明确其在风险-收益空间中的位置
    \item \textbf{慎选工具}：不是所有问题都需要深度学习，传统方法可能更安全
    \item \textbf{了解法规}：关注 EU AI Act 等监管框架，确保合规
    \item \textbf{系统化测试}：使用评估工具（如 OPIC）进行持续、一致的测试
    \item \textbf{代码优于提示}：安全约束应尽量通过代码硬规则实现，而非仅依赖 Prompt
    \item \textbf{假设会出错}：设计防御性架构，添加输出过滤层
    \item \textbf{测试偏见}：主动测试系统在不同人口统计群体中的表现差异
    \item \textbf{保持控制}：确保系统有``大红色停止按钮''
    \item \textbf{承担责任}：你的 AI，你的责任——没有例外
\end{enumerate}

\subsection{延伸阅读}

\begin{itemize}
    \item \textbf{EU AI Act 全文}：\url{https://artificialintelligenceact.eu/}
    \item \textbf{AI 希波克拉底誓言论文}：AI Magazine 上关于 AI 研究者伦理誓言的讨论
    \item \textbf{OPIC}：Comet ML 开源的 LLM 评估工具，\url{https://github.com/comet-ml/opik}
    \item \textbf{COMPAS 累犯率预测偏见}：ProPublica 对 COMPAS 系统种族偏见的调查报告
    \item \textbf{Prompt Injection 研究}：关于 LLM 提示注入攻击与防御的最新研究综述
    \item \textbf{AI Slop 现象}：关于 AI 生成内容污染互联网训练数据的讨论
    \item \textbf{Deceptive Alignment}：前沿模型在对齐评估中的欺骗行为研究
\end{itemize}

%% --- Fin del contenido --- %%

\end{document}
